%\usepackage [ISO-8859-1]{inputenc}
\input{Olympiad}

\thispagestyle{fancy}
\lstset{basicstyle=\ttfamily,
  breaklines=true}
% ----------------------------------------------------

\begin{center}
\Huge{Programmeringsolympiaden 2025}
\end{center}
\vspace{-1.2cm} 
\specialsection*{Competition rules for skolkvalet}
\begin{lista}
\item The competition takes place on a date determined by the school under {\bf four hours.
    No extra time is given for lunch or breaks.} Before the competition starts, the student
    shall decide together with the teacher whether they will use their own or a school computer.
    In both of these cases, the student must write skolkvalet in an agreed-upon location
    in the school.
    
\item The competition consists of seven tasks that shall each be solved by using a computer
  program in your programming language of choice.
\item {\bf Input can be read into the program in any way}, e.g., through the program interacting
with the user (as in the sample runs in the tasks), being entered into a graphical interface,
or data files being sent to {\em standard input}. Just agree with your teacher on how the program will be tested.
\item Your solutions will be tested with prepared input data. Each task is typically tested
with 5 test cases, each awarding 1 point if your program outputs the correct answer within
an execution time limit of {\bf 3 seconds}. No input validation is required, as it adheres
to the specifications in the task.
\item There are often different constraints for the various test cases, such as the size of
the input or other limitations. These are specified in the task. {\bf Note that the difficulty
level of a task can vary significantly depending on these differences. Therefore, it might
be easier to score partial points on a seemingly difficult task than to score full points
on a seemingly easier task.} The information about partial scoring is thus extremely
important for planning your competition.
\item The grading is performed on the same or an equivalent computer as where the code was written.
Changes to the source code are not allowed after the competition. If the program cannot be
compiled, the task will be awarded 0 points.

\item If any of the following happens, 0 points are awarded on the particular {\em testcase} being tested,
  but the program shall still be tested on the rest of the testcases.
\begin{lista}
\item The execution time of the program exceeds 3 seconds \vspace{-0.2cm}
\item Execution error (run time error) \vspace{-0.2cm}
\item Wrong answer \vspace{-0.2cm}
\end{lista}

\item Participation is individual, which means, among other things, that no exchange
of ideas or files is allowed during the competition.  
\item Permitted resources: Any written material, material installed on the computer,
and material available on the internet. 
It is \ku{not} allowed to actively communicate on the internet (e.g., chatting or asking questions in a forum);
only searching for information is permitted.  
{\bf It is also not allowed to use generative AI such as ChatGPT or GitHub Copilot.}  
Calculators are allowed.  
\item Submissions must be made in the form of source code files (e.g., uppg1...uppg7 with appropriate
file extensions) uploaded to the agreed location. Other files will not be considered. Make sure to
submit the correct version of your program.  

\end{lista}

%De högst placerade i kvalet går vidare till finalen i Göteborg där landslagsplatser till BOI i Polen och IOI i Bolivia står på spel.
\begin{center}
\Large Good luck!!
\end{center}
%\end{document}

\newpage
\specialsection*{Task 1 -- The Triangle Factory}

Tristian works at a triangle factory.  
His job is to classify different types of triangles produced in the factory.  
Tristian now asks you to write a program that can replace him.  

Tristian will give you three positive integers, $a$, $b$, and $c$.  
These are the three angles of a triangle, given in degrees.  

Your program should do the following:  
\begin{itemize}
  \item If the triangle is obtuse, the program should print ``Trubbig Triangel''.
  \item If the triangle is acute, the program should print ``Spetsig Triangel''.
  \item If the triangle is right-angled, the program should print ``Ratvinklig Triangel''.  
\end{itemize}  

To remind you of the definitions of obtuse, acute, and right-angled triangles:  
\begin{itemize}
  \item A triangle is obtuse if any of its angles are greater than $90$ degrees.  
  \item A triangle is acute if all its angles are less than $90$ degrees.  
  \item A triangle is right-angled if one of its angles is exactly $90$ degrees.  
\end{itemize}  

Your program should read the three angles $a$, $b$, $c$.  

It is guaranteed that $(1 \leq a, b, c < 180)$ holds. It is also guaranteed that the input forms a valid triangle, meaning $a + b + c = 180$.  

The program should then print whether the triangle formed by the angles $a$, $b$, and $c$ is acute, obtuse, or right-angled.  

\textbf{Scoring:} The program will be run on 5 test cases, each worth $1$ point.

\vspace{1cm}
\fe{Sample run 1}
\begin{verbatim}
Angle 1 ? 100
Angle 2 ? 40
Angle 3 ? 40

Answer: Trubbig Triangel
\end{verbatim}
The next page contains more sample runs.

\newpage

\fe{Sample run 2}
\begin{verbatim}
Angle 1 ? 60
Angle 2 ? 40
Angle 3 ? 80

Answer: Spetsig Triangel
\end{verbatim}

\fe{Sample run 3}
\begin{verbatim}
Angle 1 ? 30
Angle 2 ? 60
Angle 3 ? 90

Answer: Ratvinklig Triangel
\end{verbatim}


\newpage
\specialsection*{Task 2 -- Truls' Troubles}


%\begin{center}
%  \includegraphics[width=12cm]{sifferkrypto.pdf} \\ ~\\
%  \ku{Exempel på missuppfattning som kan ske med Alice nya chiffer.}
%\end{center}

Truls and Henry are playing table tennis against each other.

Both Truls and Henry strongly dislike math.  
More than anything, they hate calculating how many points each of them has during the game.  
To simplify scoring, they have come up with their own solution:

Instead of constantly keeping track of their points, they write a ``\texttt{T}'' if Truls
scores a point and an ``\texttt{H}'' if Henry scores a point, immediately after each won rally.

When Truls and Henry play table tennis, the scoring system works according to these three rules:  
\begin{enumerate}
  \item One point is awarded after each rally.  
  \item A player wins a match if both of the following conditions are met:  
  \begin{itemize}
    \item The player leads by at least 2 points over their opponent.  
    \item The player has at least 11 points.  
  \end{itemize}  
  For example, if Truls has 11 points and Henry has 9 points, Truls is declared the winner.
  But if Truls has 11 points and Henry has 10 points, there is no winner yet.  

  \item When a player wins according to the criteria above, the points for both players are reset to zero, and a new match begins.  
\end{enumerate}  

Now Truls and Henry have been playing for a long time, and both want to know the current score.  
They ask you to write a program that calculates the current score for them based on their notes.

Your program should first read the score string, which consists of multiple letters on a single line. The number of letters is between 1 and 30.  
Each letter will be either ``\texttt{T}'' or ``\texttt{H}''.

The program should then output the score of the ongoing match as ``\textit{$T$-$H$}'', where $T$ is
the number of points Truls has, and $H$ is the number of points Henry has.

\textbf{Scoring:}\\  
The five test cases, each worth $1$ point, have the following constraints:  
\begin{itemize}
  \item[(1-2)] $\;\;\;\;$A match always ends as soon as one player reaches 11 points.  
  \item[(3-5)] $\;\;\;\;$No additional constraints.  
\end{itemize}

\vspace{1cm}

\fe{Sample run 1}
\begin{verbatim}
The notes ? TTTHHHTHH

Answer: 4-5
\end{verbatim}

The next page contains more sample runs.


\newpage

\fe{Sample run 2}
\begin{verbatim}
The notes ? TTTTTTTTTTT

Answer: 0-0
\end{verbatim}

\fe{Explanation:} Since Truls won the match when the score reached \texttt{11-0}, the score was reset to \texttt{0-0}.

\vspace{0.5cm}

\fe{Sample run 3}
\begin{verbatim}
The notes ? THTHTHTHTHTHTHTHTHTHTHTH

Answer: 12-12
\end{verbatim}
\fe{Explanation:} Even though both have more than 11 points, neither leads by 2 points. Therefore, no winner has been declared yet.


\fe{Sample run 4}
\begin{verbatim}
The notes ? THTHTHTHTHTHTHTHTHTHTHTHHHH

Answer: 0-1
\end{verbatim}
\fe{Explanation:} This game is a continuation of sample run 3, where Henry wins the match after scoring two more points.  
The score is reset, and Henry then wins another point. Therefore, the score is \texttt{0-1}.


\newpage
\specialsection*{Task 3 -- Virus}

Boschua isn't very good at cybersecurity. Last month, he turned off the firewall
because it wouldn't let him download cat pictures. Unfortunately, this led to him getting a virus.  
The virus did something very strange: it added a lot of letters to all filenames.  

For example, it might have changed the filename \texttt{katt.png} to \texttt{k\textbf{atteg}att\textbf{o}.p\textbf{i}ng}.  
The virus can add letters but never remove them.  

Now he is looking for the file named $F$, but he has a huge number of files. Can you write a
program that reads the filename $F$ and the name of a file $H$ on his computer and determines
whether $F$ could have been transformed into $H$ by the virus?  
In other words, is it possible that the virus took the filename $F$ and added letters such that the file is now named $H$?  

Write a program that reads $F$ and $H$. Let $|F|$ be the number of characters in $F$, and the same for $|H|$.  
It holds that $1 \leq |F| \leq |H| \leq 32$. $F$ and $H$ contain only letters \texttt{a}-\texttt{z} and periods.  
It is guaranteed that $F$ contains exactly one period. The virus can add periods and any letter from \texttt{a}-\texttt{z}.  

Your program should then print ``Ja'' if $H$ could be the file Boschua is looking for, otherwise ``Nej''.  

\vspace{0.5cm}  
\textbf{Scoring:}\\  
This problem has five \textit{pairs} of test cases. In each pair, the answer to one test case is Ja and
the answer to the other is Nej.  
This ensures you do not score points if your program always prints ``Ja'' or always prints ``Nej''.  

The five pairs of test cases, each worth $1$ point, have the following constraints in addition to those already mentioned:  
\begin{itemize}  
    \item[(Pair 1)] $\;\;\;\;$Each character (and period) appears at most once in each string.  
    \item[(Pair 2)] $\;\;\;\;|H| = |F|+1$  
    \item[(Pair 3)] $\;\;\;\;|H| = |F|+2$  
    \item[(Pairs 4-5)] $\;\;\;\;$No additional constraints.  
\end{itemize} 

\vspace{0.5cm}
\fe{Sample run 1}
\begin{verbatim}
F ? katt.png
H ? kattegatto.ping

Answer: Ja
\end{verbatim}
\fe{Explanation:} The virus could have made the following changes to change $F$ to $H$:
\texttt{katt.png} $\Rightarrow$ \texttt{ka\textbf{ttega}tt\textbf{o}.p\textbf{i}ng}.

The next page contains more sample runs.

\newpage

\fe{Sample run 2}
\begin{verbatim}
F ? antivirus.exe
H ? aaantivisiiruis.egxe

Answer: Ja
\end{verbatim}
\fe{Explanation:} The virus could have made the following changes to change $F$ to $H$: \\
\texttt{antivirus.exe} $\Rightarrow$ \texttt{a\textbf{aa}ntiv\textbf{is}i\textbf{i}ru\textbf{i}s.e\textbf{g}xe}.

\fe{Sample run 3}
\begin{verbatim}
F ? kattenmilly.png
H ? kaaatttenmilly.pgn

Answer: Nej
\end{verbatim}
\fe{Explanation:} $F$ and $H$ differ by \texttt{png} and \texttt{pgn}. Because the virus cannot change the order
of letters, it could not have changed $F$ to $H$. 



\newpage
\specialsection*{Task 4 -- Siffrid's Digit Sum}

Siffrid loves playing with numbers! Right now, she is playing with the integer $N$.  

She is wondering how to create the smallest and largest numbers that have the same 
digit sum and the same number of digits as $N$.  
Can you help her?  

The digit sum of a number is defined as the sum of all its digits.  

For example, the digit sum of the number $1234$ is:  
$$
1 + 2 + 3 + 4 = 10,
$$  

and the digit sum of the number $220$ is:  
$$
2 + 2 + 0 = 4.
$$  

Write a program that first reads the integer $N$ ($1 \leq N \leq 10^9$), the number Siffrid is playing with.  

The program should then output the smallest and the largest numbers that have the same digit sum 
and the same number of digits as the number $N$.  

\vspace{0.5cm}  
\fe{Scoring:}\\  
The five test cases, each worth $1$ point, have the following additional constraints beyond those already mentioned:  

\begin{itemize}
    \item[(1-2)] $\;\;\;\;$ $N \leq 10^5$    
    \item[(3-5)] $\;\;\;\;$ No additional constraints.
\end{itemize}

\vspace{0.3cm}  
\fe{Sample run 1}\\
\vspace{-0.3cm}
\begin{verbatim}
N ? 101

Answer: 101 200
\end{verbatim}
\fe{Explanation:} We cannot get a number smaller than 101. 
Numbers like 11 and 2 have the same digit sum as $N$ and are smaller than 101, 
but they do not have the same number of digits as $N$.


\vspace{0.2cm}
\fe{Sample run 2}\\
\vspace{-0.3cm}
\begin{verbatim}
N ? 878

Answer: 599 995
\end{verbatim}
\vspace{0.2cm}

\fe{Sample run 3}\\
\vspace{-0.3cm}
\begin{verbatim}
N ? 1337

Answer: 1049 9500
\end{verbatim}


\newpage
\specialsection*{Task 5 -- The Calculator}

Rumors are spreading that the school's calculator can be hacked using a secret code.  
You've heard that if you can make the calculator display the secret Kålnami code and then press the equals button, 
it will unlock several secret mathematical functions.  

Your friend has figured out what the Kålnami code is but tells you that you can't just enter it in any way you like.  
Initially, the calculator displays $0$. After that, you are allowed to perform the following operations:  
\begin{itemize}
  \item Multiply the number displayed by $M$.
  \item Add any number $k$ to the number displayed, where $0 \leq k \leq M-1$.
\end{itemize}  

Fortunately, your friend has also figured out what $M$ is.  
This complicated process will only work if you use the minimum number of operations possible.  
You now want to determine how many operations are required, given the Kålnami code and the number $M$.  

Write a program that first reads the integer $N$ ($1 \leq N \leq 10^9$), representing the Kålnami code.  
Then it should read the digit $M$ ($2 \leq M \leq 9$), as described above.  

The program should then output a single integer: the minimum number of operations required to make the calculator display the Kålnami code, $N$.  

\fe{Scoring:}\\  
The five test cases, each worth $1$ point, have the following additional constraints beyond those already mentioned:  
\begin{itemize}
    \item[(1)] $\;\;\;\;M=3, N \leq 10$ 
    \item[(2)] $\;\;\;\;M=2$ 
    \item[(3)] $\;\;\;\;N \leq 10^5$
    \item[(4-5)] $\;\;\;\; $No additional constraints.
\end{itemize}  

\fe{Sample run 1}  
\begin{verbatim}
N ? 4
M ? 2

Answer: 3
\end{verbatim}  

\fe{Explanation:}  
One way to enter the code using the minimum number of operations is as follows:  
\begin{enumerate}
  \item $+1$
  \item $\times 2$
  \item $\times 2$
\end{enumerate}  

See the next page for another example run.  

\newpage  

\fe{Sample run 2}  
\begin{verbatim}
N ? 33
M ? 3

Answer: 5
\end{verbatim}  

\fe{Explanation:}  
An optimal solution is to perform the following operations:  
\begin{enumerate}
  \item $+1$
  \item $\times 3$
  \item $\times 3$
  \item $+2$
  \item $\times 3$
\end{enumerate}  



\newpage
\specialsection*{Task 6 -- The Wizard Theodor}

A smaller army of $N$ monsters is just outside Tästerås! 
Now it's up to Trollkarlen Theodor to defeat them and save the city. 
Each monster has a certain amount of health. To damage the monsters, 
Theodor can fire magical explosions. Each time he fires an explosion, 
he targets a specific monster. Then the following happens:  
\begin{itemize}
  \item All monsters lose $A$ health from the explosion.
  \item The monster Theodor targets loses an additional $S$ health because it is at the center of the explosion. Thus, the targeted monster loses $S+A$ health in total from that explosion.
\end{itemize}
He can freely choose which monster to target each time he fires an explosion.  
Since firing explosions is exhausting, he wants to know the minimum number of explosions he needs to fire to defeat all the monsters if he targets optimally.  

Write a program that first reads the integers $N$ ($1 \leq N \leq 10$), $S$ ($1 \leq S \leq 10^9$), and $A$ ($0 \leq A \leq 10^9$), as described above. 
Then, the program should read $N$ integers $h_1, h_2, \ldots, h_N$ ($1 \leq h_i \leq 10^9$), representing the health of each monster.  

The program should output the minimum number of explosions Theodor needs to fire to defeat all the monsters. A monster is defeated if its health is reduced to $0$ or less.  

\fe{Scoring:}\\  
The five test cases, each worth $1$ point, have the following additional constraints beyond those already mentioned:

\begin{itemize}
    \item[(1)] $\;\;\;\; A=0$
    \item[(2-3)] $\;\;\;\; h_i \leq 100$
    \item[(4-5)] $\;\;\;\; $No additional constraints.
\end{itemize}


\fe{Sample run 1}
\begin{verbatim}
N, S, A ? 3 2 1
Monsters' health ? 7 2 3

Answer: 3
\end{verbatim}
\fe{Explanation:} An optimal way to fire explosions is as follows:
\vspace{-0.2cm}
\begin{enumerate}
  \item Fire an explosion at the monster with 7 health. The monsters now have \verb+4,1,2+ health remaining.
  \item Fire an explosion at the monster with 4 health. The monster with 1 health is defeated. The remaining monsters have \verb+1,1+ health left.
  \item Fire an explosion at either of the remaining monsters. All monsters are now defeated, and Theodor is finished.
\end{enumerate}

\vspace{-0.2cm}
\fe{Sample run 2}
\begin{verbatim}
N, S, A ? 4 3 2
Monsters' health ? 7 8 9 10

Answer: 4
\end{verbatim}



\specialsection*{Task 7 -- Theodor, The Wizard}


This problem is based on the problem \textit{The Wizard Theodor}.  

Evil Olle, the mastermind behind all evil, is tired of Theodor always ruining his plans. 
Now, Olle has come up with a new devious scheme, and to carry it out, he needs to keep Theodor occupied. 
At his disposal, he has $N$ monsters, and since he is immensely powerful, he can assign these monsters different amounts of health. 
Olle now wants to count the number of ways he can distribute different health values among the monsters such that 
Theodor needs to fire exactly $K$ magical explosions to defeat all the monsters. 
The only requirement for a monster's health is that it must be greater than $0$.  

Just like in the problem \textit{The Wizard Theodor}, 
Theodor damages monsters by targeting a specific one and firing an explosion. 
The targeted monster loses $S$ health, and then all monsters (including the targeted one) lose $A$ health each. 
A monster is defeated if it has $0$ or less health. Theodor targets monsters in such a way as to minimize the number of explosions he needs to fire.  

Write a program that reads the integers $N, K, S, A$ ($1 \leq N, K, S \leq 2\cdot10^5$, $0 \leq A \leq 2\cdot10^5$), which are the four numbers described above.  

The program should then output the number of ways Olle can assign health values to the 
monsters such that Theodor needs to fire exactly $K$ explosions to defeat all of them. 
Since this number can be very large, it should be output \textit{modulo} $10^9+7$. 
If you only want to solve the first test case, you do not need to worry about modulo as the answer will be small.



\fe{Scoring:}\\
The five test cases, each worth $1$ point, have the following additional constraints beyond those already mentioned:
\begin{itemize}
    \item[(1)] $\;\;\;\; A = 0$, $N, K, S \leq 5$
    \item[(2)] $\;\;\;\; A = 0$, $S \leq 10$, $N, K \leq 70$
    \item[(3)] $\;\;\;\; N, K, S, A \leq 150$
    \item[(4)] $\;\;\;\; N, K, S, A \leq 1000$
    \item[(5)] $\;\;\;\;$No additional constraints.
\end{itemize}


\fe{Sample run 1}
\begin{verbatim}
N, K, S, A ? 1 1 5 0

Answer: 5
\end{verbatim}
\fe{Explanation:} There is one monster, and the possible health values it can have such that Theodor defeats it with exactly one explosion are \verb+1, 2, 3, 4, 5+.

See the next page for more sample runs.

\newpage

\fe{Sample run 2}
\begin{verbatim}
N, K, S, A ? 2 2 1 1

Answer: 10
\end{verbatim}
\fe{Explanation:} The 10 possible distributions of health values for the monsters such that Theodor defeats both with exactly two explosions are:
\begin{verbatim}
  1 3
  1 4
  2 2
  2 3
  2 4
  3 1
  3 2
  3 3
  4 1
  4 2
\end{verbatim}

\vspace{1cm}
\fe{Sample run 3}
\begin{verbatim}
N, K, S, A ? 10 20 30 40

Answer: 435462573
\end{verbatim}

\end{document}
